\section{Matroid union and Edmonds' partition theorem}
\label{sec:matroid-union}

This chapter formalizes the classical \textbf{matroid-union theorem}
of Nash-Williams~\cite{nashWilliams1966} and Edmonds, together with
the matroid-from-submodular-function construction and polymatroid
rank formula of Edmonds~\cite{edmonds1970} on which it rests, and
\textbf{Edmonds' matroid-partition theorem}~\cite{edmonds1965}. These
are the abstract-matroid prerequisites of the body-bar route
(\cref{sec:body-bar}): the $k$-frame matroid is identified with a
$k$-fold matroid union, and the count characterization underlying
Tutte--Nash-Williams tree-packing follows from the partition rank
formula. For a textbook treatment see Oxley~\cite{oxley2011}.

\paragraph{Provenance.}
The Lean development of this chapter is ported from Peter Nelson's
\href{https://github.com/apnelson1/Matroid}{\texttt{apnelson1/Matroid}}
package (the same package that supplies \texttt{Matroid.ofFun} in
\cref{sec:rigidity-matroid} and \texttt{Graph.cycleMatroid} in
\cref{sec:body-bar}), whose shelved development already contains
complete proofs of all the results below. Those proofs rest on a
matroid constructor that has since been superseded upstream, so this
project vendors them into \texttt{CombinatorialRigidity/Matroid/},
rebased onto the package's current constructor and retaining Nelson's
authorship. The package is Apache-2.0 licensed (as are this project
and mathlib); the per-file headers under
\texttt{CombinatorialRigidity/Matroid/} record the attribution and the
list of modifications.

\subsection{Matroids from submodular functions}
\label{sec:matroid-union-submodular}

The union theorem realizes a matroid union as the matroid induced by a
submodular function, so we begin with that construction and its rank
formula, following Edmonds~\cite{edmonds1970}. A non-decreasing
submodular function on a finite ground set induces a matroid; on a
polymatroid rank function the rank of any set is given by a
minimization formula.

\begin{definition}[Submodular function]
  \label{def:submodular}
  \lean{Matroid.Submodular}
  \leanok
  An integer-valued set function $f$ on the subsets of a finite
  ground set $E$ is \emph{submodular} if
  $f(X \cap Y) + f(X \cup Y) \le f(X) + f(Y)$ for all $X, Y$.
\end{definition}

\begin{definition}[Polymatroid rank function]
  \label{def:polymatroidFn}
  \lean{Matroid.PolymatroidFn}
  \leanok
  \uses{def:submodular}
  A \emph{polymatroid (rank) function} on $E$ is a non-decreasing
  submodular function $f$ with $f(\emptyset) = 0$
  (Edmonds~\cite{edmonds1970}).
\end{definition}

\begin{definition}[Matroid of a submodular function]
  \label{def:ofSubmodular}
  \lean{Matroid.ofSubmodular, Matroid.circuit_ofSubmodular_iff,
    Matroid.indep_ofSubmodular_iff}
  \leanok
  \uses{def:submodular}
  For a non-decreasing submodular $f$, the \emph{induced matroid}
  $M_f$ on $E$ has as its circuits the minimal non-empty sets $C$
  with $f(C) < |C|$; equivalently, a set $I$ is independent in $M_f$
  if and only if $|I'| \le f(I')$ for every non-empty $I' \subseteq I$
  (Edmonds~\cite{edmonds1970}).
\end{definition}

\begin{definition}[Matroid of a polymatroid function]
  \label{def:ofPolymatroidFn}
  \lean{Matroid.ofPolymatroidFn, Matroid.indep_ofPolymatroidFn_iff}
  \leanok
  \uses{def:ofSubmodular, def:polymatroidFn}
  The matroid $M_f$ of \cref{def:ofSubmodular} specialized to a
  polymatroid rank function $f$.
\end{definition}

\begin{lemma}[Polymatroid rank formula]
  \label{lem:polymatroid-rank}
  \lean{Matroid.polymatroid_rank_eq}
  \leanok
  \uses{def:ofPolymatroidFn}
  For a polymatroid rank function $f$ and any $X \subseteq E$, the
  rank of $X$ in $M_f$ is
  \[
    r_{M_f}(X) \;=\; \min_{Y \subseteq X}\bigl(f(Y) + |X \setminus Y|\bigr).
  \]
\end{lemma}
\begin{proof}
  \leanok
  \uses{def:ofSubmodular}
  Edmonds~\cite{edmonds1970}. The lower bound holds since any
  independent subset of $X$ meets each $Y$ in at most $f(Y)$ elements
  and $X \setminus Y$ in at most $|X \setminus Y|$; a basis of $X$
  attains the minimum at the $Y$ on which it is ``tight''.
\end{proof}

\subsection{Matroid union}
\label{sec:matroid-union-union}

The union of a finite family of matroids on a common ground set is
again a matroid --- the Nash-Williams--Edmonds union
theorem~\cite{nashWilliams1966}. It is realized in two ways: as the
transversal matroid of a bipartite matching between $E$ and the direct
sum of the $M_{j}$, and as the induced matroid $M_f$ of the submodular
function $f(X) = \sum_{j} r_{M_{j}}(X)$ (\cref{def:ofPolymatroidFn}).
The first realization gives the decomposition characterization of the
union's independent sets; the second gives Edmonds' matroid-partition
rank formula.

\begin{definition}[Matroid union of an indexed family]
  \label{def:matroid-union}
  \lean{Matroid.Union, Matroid.union}
  \leanok
  \uses{def:ofPolymatroidFn, lem:polymatroid-rank}
  Given a finite indexed family of matroids $M_{1}, \dots, M_{k}$ on a
  common ground set $E$, their \emph{union} $\bigcup_{j} M_{j}$ is the
  matroid on $E$ whose independent sets are the unions
  $I_{1} \cup \dots \cup I_{k}$ of sets $I_{j}$ independent in $M_{j}$.
\end{definition}

\begin{lemma}[Union independence by decomposition]
  \label{lem:union-indep-iff}
  \lean{Matroid.union_indep_iff, Matroid.union_indep_iff'}
  \leanok
  \uses{def:matroid-union}
  A set $I$ is independent in $\bigcup_{j} M_{j}$ if and only if it
  decomposes as $I = \bigcup_{j} I_{j}$ with each $I_{j}$ independent
  in $M_{j}$.
\end{lemma}
\begin{proof}
  \leanok
  The transversal realization (\cref{def:matroid-union}) matches the
  ground set $E$ into the direct sum of the $M_{j}$, and a set is
  independent exactly when it admits such a matching. Reading off the
  factor of each matched element gives a decomposition
  $I = \bigcup_{j} I_{j}$ with $I_{j}$ independent in $M_{j}$, and
  conversely any such decomposition matches; discarding elements shared
  between factors makes the $I_{j}$ pairwise disjoint.
\end{proof}

\begin{lemma}[Rado's theorem]
  \label{lem:rado}
  \lean{Matroid.rado}
  \leanok
  \uses{lem:polymatroid-rank}
  Let $M$ be a matroid on a ground set $\alpha$, and
  $(A_i)_{i \in \iota}$ a finite family of finite subsets of $\alpha$.
  Some injective transversal $i \mapsto e_i$ with $e_i \in A_i$ for
  every $i$ has range independent in $M$ if and only if
  $|K| \le r_M\bigl(\bigcup_{i \in K} A_i\bigr)$ for every
  $K \subseteq \iota$.
\end{lemma}
\begin{proof}
  \leanok
  \uses{lem:polymatroid-rank}
  This is Rado's theorem~\cite[Thm 11.2.2]{oxley2011} (Rado 1942). Apply
  the submodular generalization of Hall's marriage condition to the rank
  function $f(S) = r_M(S)$, which is monotone and submodular; the
  representative-existence half of that condition yields the independent
  transversal, while the rank-condition side is exactly the hypothesis.
\end{proof}

\begin{theorem}[Edmonds' matroid-partition rank formula]
  \label{thm:matroid-partition-rank}
  \lean{Matroid.Union_rank_eq, Matroid.Union_pow_rk_eq}
  \leanok
  \uses{def:matroid-union, lem:union-indep-iff, lem:polymatroid-rank}
  Let $M_{1}, \dots, M_{k}$ be matroids on a common finite ground set
  $E$, and $E' \subseteq E$. The rank of $E'$ in the union
  $\bigcup_{j} M_{j}$ equals
  $\min_{X \subseteq E'} \bigl(|E' \setminus X| + \sum_{j}
  r_{M_{j}}(X)\bigr)$. In particular, $E'$ is independent in the union
  if and only if $|E''| \le \sum_{j} r_{M_{j}}(E'')$ for every
  $E'' \subseteq E'$.
\end{theorem}
\begin{proof}
  \leanok
  \uses{lem:polymatroid-rank, lem:rado}
  Apply the polymatroid rank formula (\cref{lem:polymatroid-rank}) to
  $f(X) = \sum_{j} r_{M_{j}}(X)$, which realizes the union
  (\cref{def:matroid-union}); the matching that exhibits the union's
  independent sets is furnished by Rado's theorem (\cref{lem:rado}).
  This is Edmonds' matroid-partition theorem~\cite{edmonds1965}.
\end{proof}

Specializing this formula is the entry point for the body-bar route:
taking the $M_{j}$ to be $k$ copies of a graphic matroid turns the
independence count into $(k, k)$-sparsity, the matroid-side input to
the Tutte--Nash-Williams tree-packing count (\cref{sec:body-bar}).
